\section{Introduction}
Vector search, the task of finding the $k$-nearest neighbors (KNN) of a query vector against a database of high-dimensional vectors, is fundamental to modern machine learning based search~\cite{grbovic2018real, haldar2019applying, hashemi2021neural, qin2021mixer} and recommendation systems~\cite{okura2017embedding, liu2017related, wang2018billion, pal2020pinnersage, liu2022monolith}. In these applications, a vector represents an item in a metric space, and the distance between vectors reflects semantic similarity. However, performing exact KNN search becomes computationally infeasible on large datasets due to the high dimensionality and volume of data.

To address this challenge, practitioners use approximate nearest neighbor (ANN) indexes, which trade off a controlled amount of search accuracy (recall) for significant reductions in latency. Among these, two broad classes dominate in practice: graph-based and partitioned indexes, each with distinct performance characteristics under dynamic workloads.


% Among these indexes, there are two widely-used types: graph-based indexes and partitioned indexes.

Maintaining low latency, high recall vector search under \textbf{dynamic and skewed workloads} remains a significant challenge for existing indexes. Real-world applications often exhibit non-uniform query distributions and evolving data. For example, in an example Wikipedia search application, popular pages like \textit{Lionel Messi} or \textit{LeBron James} receive disproportionately more queries, resulting in \emph{skewed read patterns}. Additionally, pages are frequently added, updated, or deleted, causing \emph{skewed update patterns} that change over time ~\cite{baranchuk_dedrift_2023}. These factors degrade the performance of existing indexes, leading to increased query latency and reduced recall.

\emph{Graph-based indexes}, such as HNSW~\cite{10.1109/TPAMI.2018.2889473}, DiskANN ~\cite{10.5555/3454287.3455520, singh2021freshdiskann}, and SVS ~\cite{aguerrebere2024locally} construct a proximity graph where each node (vector) is connected to its approximate neighbors. Queries traverse the graph to find approximate nearest neighbors, typically achieving high recall with low latency. However, these indexes face challenges with dynamic workloads because updating the graph structure to accommodate frequent insertions and deletions is computationally intensive~\cite{xu_spfresh_2023}, due to the random access patterns involved in graph traversal and modification.

\emph{Partitioned indexes}, such as SCANN \cite{guo_accelerating_2020, soar_2023}, SPANN \cite{xu_spfresh_2023, chen_spann_nodate}, and  Faiss-IVF \cite{douze2024faisslibrary}, partition the vectors using a clustering algorithm (e.g k-means). Queries are processed by scanning a subset of partitions, balancing recall and latency by adjusting the number of partitions scanned (denoted as \emph{nprobe}).  While attractive due to their simplicity, partitioned indexes face a significant search latency gap when compared with graph indexes. For example, on the \textsc{MSTuring10M} benchmark \cite{bigann}, we found Faiss-IVF takes 44ms per search query while Faiss-HNSW takes only 6.8ms. 
On the other hand, supporting updates in partitioned indexes is less expensive than for graph indexes, as the index structure needs minimal modification when adding or removing vectors. But, existing approaches struggle with dynamic and skewed workloads because they do not adapt to changing access patterns, leading to \emph{imbalanced} partitions that degrade query latency. Recent work has been proposed to resolve imbalances in dynamic workloads by splitting and reclustering imbalanced partitions \cite{xu_spfresh_2023, baranchuk_dedrift_2023}, however, we find these methods degrade recall as nprobe needs to change as the index structure changes.

In this work, we study the problem of minimizing query latency to meet a fixed recall target for dynamic vector search workloads, where both the queries and the base vectors can change over time. To address this problem, we develop \textbf{Quake}, a partitioned index for ANN search that minimizes query latency by adapting the index structure to the workload. \edit{Quake’s two primary algorithmic contributions are:}


First, Quake employs an \textbf{adaptive hierarchical partitioning} scheme that modifies the partitioning by minimizing the cost (derived from a proposed \textbf{cost model}) of a query. The cost model tracks partition sizes and access frequencies as the workload is processed and determines which partitions are most negatively contributing to overall query latency. Once identified, we split or merge these partitions based on expected cost reduction derived from our proposed cost model. We also demonstrate that our maintenance procedure is stable and converges to a local minimum of the cost model.

Second, we design an \textbf{adaptive partition scanning} scheme that adjusts the number of partitions scanned on-the-fly to meet recall targets for individual queries. We do this by maintaining a recall estimate during query processing based on A) the geometry of the partitioning and B) intermediate results of the query, and once the estimate exceeds the recall target, query processing terminates and the results are returned.

\edit{Furthermore, Quake utilizes NUMA-aware parallelism to maximize memory bandwidth usage on multi-core machines.}

It is a significant challenge to evaluate indexing approaches due to the lack of availability of benchmarks for online vector search. To address this challenge and comprehensively evaluate our approach, we A) prepare a \textbf{Wikipedia vector search workload} derived from publicly available query and update patterns of Wikipedia pages and B) develop a \textbf{workload generator} for creating workloads with configurable query and update patterns. We will publicly release the Wikipedia workload and workload generator as evaluation tools for the community to use. Using these, we conduct a comprehensive evaluation of Quake in comparison to seven baseline approaches.

\begin{enumerate}[leftmargin=*, itemsep=-1pt, topsep=0pt]
    \item Quake achieves the lowest search time across all dynamic workloads compared to state-of-the-art graph indexes, with $1.5$-$13\times$ lower search latency than HNSW, DiskANN, and SVS while having $18$-$126\times$ lower update latency.
    \item We also find that APS matches the nprobe of an oracle across recall targets on \textsc{Sift1M}, with only a $17$-$29\%$ increase in latency relative to the oracle. 
    \item APS performs on-par or better than existing early termination methods~\cite{zhang2023fast, chen_spann_nodate, li2020improving} and requires no offline tuning. %Additionally, in contrast to existing methods for early termination, APS requires no offline tuning.
    \item Quake's NUMA-aware query processing exhibits linear scalability and high memory bandwidth utilization on the \textsc{MSTuring100M} dataset. Quake achieves $20\times$ and $4\times$ lower query latency compared to single-threaded and non-NUMA aware configurations, respectively.% version and $4\times$ lower latency compared to a non-NUMA aware configuration.
\end{enumerate}